\begin{landscape}
\begingroup
\footnotesize
\setlength{\tabcolsep}{4pt}
\renewcommand{\arraystretch}{1.18}

\begin{longtable}{|>{\raggedright\arraybackslash}p{0.18\linewidth}|>{\raggedright\arraybackslash}p{0.20\linewidth}|>{\raggedright\arraybackslash}p{0.25\linewidth}|>{\raggedright\arraybackslash}p{0.29\linewidth}|}
\caption{Đề xuất xử lý các tình huống chưa phù hợp trong quá trình kiểm nghiệm}\label{tab:de-xuat-xu-ly-ket-qua} \\
\hline
\textbf{Tình huống} & \textbf{Biểu hiện} & \textbf{Nguyên nhân có thể} & \textbf{Đề xuất xử lý} \\
\hline
\endfirsthead

\multicolumn{4}{c}{\tablename~\thetable{} -- tiếp theo} \\
\hline
\textbf{Tình huống} & \textbf{Biểu hiện} & \textbf{Nguyên nhân có thể} & \textbf{Đề xuất xử lý} \\
\hline
\endhead

\hline
\multicolumn{4}{r}{\textit{Còn tiếp trang sau}} \\
\endfoot

\hline
\endlastfoot

Độ lặp không đạt yêu cầu
& Chênh lệch giữa hai lần phân tích lớn; RPD vượt giới hạn chấp nhận của phương pháp
& Mẫu chưa được nghiền và trộn đồng nhất trước khi chia mẫu phân tích
& Nghiền và đồng nhất lại mẫu; lấy lại hai phần mẫu độc lập để tiến hành phân tích song song. \\
\hline

&
& Sai số trong quá trình cân mẫu, hút pipet, định mức hoặc pha loãng
& Kiểm tra lại cân, pipet và dụng cụ định mức; thực hiện thao tác đúng quy trình và hạn chế thay đổi người thao tác giữa các lần lặp. \\
\hline

&
& Sai khác trong quá trình xử lý mẫu, thời gian phản ứng, nhiệt độ hoặc điều kiện thí nghiệm
& Kiểm soát thống nhất thời gian, nhiệt độ, lượng hóa chất và trình tự thao tác giữa các lần thử. \\
\hline

&
& Xác định điểm kết thúc chuẩn độ chưa đồng nhất
& Chuẩn độ chậm khi gần điểm cuối; sử dụng cùng loại chỉ thị và tiêu chí nhận biết màu; thực hiện mẫu trắng song song. \\
\hline

&
& Thiết bị phân tích hoạt động chưa ổn định, đặc biệt đối với máy quang kế ngọn lửa
& Kiểm tra đường chuẩn, dung dịch chuẩn kiểm soát, độ ổn định của ngọn lửa, hệ thống hút và đầu phun; vệ sinh thiết bị trước khi đo lại. \\
\hline

&
& Hóa chất hoặc dung dịch chuẩn có nồng độ không ổn định
& Kiểm tra hạn sử dụng, điều kiện bảo quản; chuẩn hóa lại dung dịch chuẩn khi cần thiết. \\
\hline

&
& --
& \textbf{Thực hiện lại toàn bộ phép thử từ bước cân mẫu. Chỉ lấy kết quả trung bình khi các lần phân tích mới đáp ứng yêu cầu độ lặp.} \\
\hline

Kết quả thành phẩm khác với chỉ tiêu công bố
& Hàm lượng N, \(\mathrm{P_2O_5}\), \(\mathrm{K_2O}\) hoặc chất hữu cơ thấp hơn giới hạn cho phép so với mức đăng ký
& Có thể xuất phát từ sai số phân tích hoặc quá trình chuẩn bị mẫu
& Trước tiên kiểm tra lại phép tính, hệ số pha loãng, đường chuẩn, mẫu trắng và điều kiện phân tích; sau đó lấy mẫu mới và phân tích lại để xác nhận. \\
\hline

&
& Thành phần dinh dưỡng thực tế của nguyên liệu khác với giá trị sử dụng khi tính công thức phối trộn
& Kiểm tra lại kết quả phân tích nguyên liệu đầu vào và hồ sơ chất lượng của từng lô nguyên liệu. \\
\hline

&
& Sai số trong quá trình định lượng nguyên liệu hoặc tỷ lệ phối trộn
& Kiểm tra phiếu phối liệu, hệ thống cân định lượng và lượng nguyên liệu thực tế đưa vào mẻ sản xuất. \\
\hline

&
& Hỗn hợp sau phối trộn chưa đồng đều
& Kiểm tra thời gian và điều kiện phối trộn; lấy mẫu tại nhiều vị trí của lô để đánh giá mức độ đồng nhất. \\
\hline

& Hàm lượng \(\mathrm{K_2O}\) hoặc \(\mathrm{P_2O_5}\) thấp bất thường
& Nguồn nguyên liệu kali hoặc lân chưa phù hợp, lượng nguyên liệu cấp vào dây chuyền sai lệch hoặc hỗn hợp phân bố không đồng đều
& Kiểm tra riêng nguồn nguyên liệu kali hoặc lân, lượng nguyên liệu được cấp vào dây chuyền và khả năng phân bố không đồng đều trong hỗn hợp. \\
\hline

& Độ ẩm thành phẩm cao hơn giới hạn
& Nguyên liệu có độ ẩm cao; sản phẩm hút ẩm trong quá trình sản xuất hoặc bảo quản; quá trình làm khô chưa phù hợp
& Kiểm tra độ ẩm nguyên liệu, điều kiện sấy hoặc làm khô, thời gian lưu sản phẩm và điều kiện kho; hạn chế tiếp xúc với môi trường có độ ẩm cao. \\
\hline

& Hàm lượng dinh dưỡng cao hơn mức đăng ký
& Có thể do sai số pha loãng, đường chuẩn hoặc lượng nguyên liệu dinh dưỡng thực tế cao hơn công thức
& Kiểm tra lại phép phân tích và công thức phối trộn; nếu kết quả phân tích lại vẫn tương tự thì đánh giá mức độ phù hợp theo giới hạn pháp lý áp dụng cho sản phẩm. \\
\hline

&
& --
& \textbf{Nếu kết quả phân tích lại vẫn không phù hợp, cần tạm thời chưa kết luận lô đạt chất lượng và chuyển thông tin cho Phòng Kỹ thuật/Bộ phận Sản xuất để xác định nguyên nhân, điều chỉnh quá trình sản xuất và kiểm nghiệm lại trước khi xuất sản phẩm.} \\

\end{longtable}
\endgroup
\end{landscape}
